\begin{homeworkProblem}[Seção 2-15]
  Sejam $u,v:U\to\mathbb{R}$ funções de classe $C^{2}$ definidas em um aberto $U\subseteq \mathbb{R}^{n}$ e $\Omega\subseteq U$ é um domínio compacto com fronteira regular e $\eta$ o campo normal unitário apontando para fora.\\
  Obtenha a Segunda Fórmula de Green:
  \begin{align*}
    \int_{\Omega}\bigl(u\Delta v-v\Delta u\bigr)=\int_{\partial\Omega}\left(u\frac{\partial v}{\partial\eta}-v\frac{\partial u}{\partial\eta}\right).
  \end{align*}
  Daí obtenha a fórmula
  \begin{align*}
    \int_{\Omega}\Delta u=\int_{\partial\Omega}\frac{\partial u}{\partial\eta}.
  \end{align*}
  \begin{solution}
    Note que
    \begin{align*}
      udv=uv_{x_1}dx^{1}+uv_{x_2}dx^{2}+\cdots+uv_{x_n}dx^{n},
    \end{align*}
    aplicando o operador de Hodge
    \begin{align*}
      \star udv=uv_{x_1}dx^{2}\wedge\cdots\wedge dx^{m}+\cdots+(-1)^{n-1}uv_{x_n}dx^{1}\wedge\cdots\wedge dx^{n-1}.
    \end{align*}
    Portanto,
    \begin{align*}
      div(udv)=d(\star udv)=&\frac{\partial uv_{x_1}}{\partial x_1}dx^{1}\wedge\cdots\wedge dx^{n}+\cdots+\frac{\partial uv_{x_n} }{\partial x_n}dx^{1}\wedge\cdots\wedge dx^{n}\\
      =&\left( \sum_{i=1}^{n}\frac{\partial uv_{x_i} }{\partial x_i} \right)dx^{1}\wedge\cdots\wedge dx^{n}\\
      =&\left( \left\langle du,dv \right\rangle+u\Delta v \right)dx^{1}\wedge\cdots\wedge dx^{n}.
    \end{align*}
    Assim, usando a simetría do produto interno e o teorema de divergencia temos que
    \begin{align*}
      \int_{\Omega}u\Delta v-v\Delta u=&\int_{\Omega}\left\langle du,dv \right\rangle+ u\Delta v-\left\langle dv,du \right\rangle-v\Delta u\\
      =&\int_{\partial \Omega}\left\langle udv,\eta \right\rangle-\left\langle vdu,\eta \right\rangle\\
      =&\int_{\partial \Omega}u\frac{\partial v}{\partial \eta}-v\frac{\partial u}{\partial \eta}.
    \end{align*}
    Agora, note que se pegamos $v=1$ e multiplicando por $-1$ temos que como $dv=0$ e $\Delta v=0$, então 
    \begin{align*}
      \int_{\Omega}\Delta u=\int_{\partial\Omega}\frac{\partial u}{\partial \eta}.
    \end{align*}
    O que nos permite concluir o exercício.
  \end{solution}
\end{homeworkProblem}
